\documentclass[11pt]{article}
\usepackage[a4paper,margin=1in]{geometry}
\usepackage{amsmath,amssymb,amsthm,booktabs,microtype}
\newtheorem{theorem}{Theorem}[section]
\newtheorem{proposition}[theorem]{Proposition}
\theoremstyle{remark}\newtheorem{remark}[theorem]{Remark}
\newcommand{\Tr}{\operatorname{Tr}}
\newcommand{\Schatten}{\mathcal S}
\title{A Modulus-Complete Spectral-Tail Certificate for the Noisy Hardy Determinant}
\author{Bin Wang}\date{July 2026}
\begin{document}\maketitle
\begin{abstract}
The fixed-rank natural-$L^2$ quotient fails at zero noise, while the
variable-rank criterion of RH-279 was previously uninstantiated.  We activate
its tail branch without constructing a Riesz projector.  After removing the
two peripheral eigenvalues of the Hardy-scaled folded Gaussian operator, take
every algebraic eigenvalue of modulus greater than $q=1/2$ as the spectral
head and put the remaining eigenvalues on a normal diagonal operator
$C_\sigma$.  The sharp Hilbert--Schmidt mass law bounds their squared sum by
$\sigma^{-1}$ for small noise.  With $R=7/5$ and
$m_\sigma=\lceil4\log(1/\sigma)\rceil$, the trace norm, operator norm,
prefix, and root-rate hypotheses of RH-279 hold, and the full logarithmic
tail beyond $m_\sigma$ vanishes uniformly on $|z|\le R$.  This proves a
projection-free spectral-tail theorem.  It does not identify the noisy head
with the monodromy counterloop.
\end{abstract}

\section{Spectral realization}
Let $A_\sigma=K_\sigma/r_H$, $r_H=0.85$, and remove the Perron and
negative-parity eigenvalues from its algebraic spectrum.  Denote the
remaining sequence, counted with algebraic multiplicity, by
$(\mu_j(\sigma))_{j\ge1}$.  RH-276 and the Hilbert--Schmidt eigenvalue
inequality give, for sufficiently small $\sigma$,
\begin{equation}\label{eq:mass}
 M_\sigma:=\sum_j|\mu_j(\sigma)|^2
 \le \|A_\sigma\|_{\Schatten_2}^2\le\sigma^{-1}.
\end{equation}
Fix $q=1/2$.  The head
$H_\sigma=\{\mu_j:|\mu_j|>q\}$ is finite.  On the Hilbert space
$\ell^2(\{j:|\mu_j|\le q\})$, define
\[
 C_\sigma=\operatorname{diag}(\mu_j:|\mu_j|\le q).
\]
The canonical-product factorization from RH-234 shows that $C_\sigma$
realizes exactly the complementary projection-free $\det_2$ factor.

\section{Variable-block theorem}
\begin{theorem}[instantiated RH-279 certificate]\label{thm:certificate}
For every integer $m\ge2$,
\begin{align}
 \|C_\sigma^m\|_1&\le M_\sigma q^{m-2},&
 \|C_\sigma^m\|&\le q^m,\\
 \|C_\sigma^r\|&\le q^r\quad(0\le r<m),&
 |\Tr C_\sigma^n|&\le M_\sigma q^{n-2}\quad(n\ge2).
\end{align}
Let $R=7/5$ and $m_\sigma=\lceil4\log(1/\sigma)\rceil$.  Then
\[
 \limsup_{\sigma\downarrow0}
 \|C_\sigma^{m_\sigma}\|_1^{1/m_\sigma}R
 \max\{1,qR\}
 \le \frac7{10}e^{1/4}<1.
\]
Consequently
\begin{equation}\label{eq:tail}
 \sum_{n\ge m_\sigma}
 \frac{|\Tr C_\sigma^n|R^n}{n}\longrightarrow0.
\end{equation}
\end{theorem}
\begin{proof}
The operator is normal and diagonal.  Since $|\mu_j|\le q$ on the tail,
\[
 \sum_j|\mu_j|^m\le q^{m-2}\sum_j|\mu_j|^2,
\]
which gives every displayed norm and trace estimate.  By
\eqref{eq:mass},
\[
 \limsup (M_\sigma q^{-2})^{1/m_\sigma}\le e^{1/4}.
\]
Because $qR=7/10$, the root-rate estimate follows.  Finally,
\[
 \sum_{n\ge m}\frac{|\Tr C_\sigma^n|R^n}{n}
 \le \frac{M_\sigma q^{-2}(qR)^m}{m(1-qR)}.
\]
For the chosen clock the right side is
$O(\sigma^{4\log(10/7)-1}/\log(1/\sigma))$, whose exponent is
$0.426699\ldots>0$.
\end{proof}

\section{Head rank and exact scope}
\begin{proposition}
The modulus-complete head obeys
$\#H_\sigma\le M_\sigma/q^2\le4\sigma^{-1}$.
\end{proposition}
\begin{proof}
Every head eigenvalue contributes more than $q^2$ to the squared spectral
mass in \eqref{eq:mass}.
\end{proof}

\begin{remark}
$C_\sigma$ is a normal spectral realization of the canonical-product tail.
It need not be similar to a bounded physical compression, and no estimate on
the ill-conditioned Riesz projectors of RH-232 follows.  The theorem supplies
the variable-rank tail obligation, not the head-to-monodromy coefficient
bridge.  Gates A--E remain false/open; no Hilbert--Polya operator, Riemann-zero
identification, zeta-divisor equality, or RH conclusion is asserted.
\end{remark}
\end{document}
